% TOP_PROBLEM: {"id": 2302082, "title": "Research Problems in Function Theory — Problem 2.82", "category_id": 8, "category": {"id": 8, "name": "analysis", "display_name": "Analysis", "description": "Limits, continuity, calculus, and function theory.", "slug": "analysis", "order_index": 8, "created_at": "2026-07-31T15:26:25.671Z"}, "status": "open", "problem_number": "AMR-022-2082", "set_id": 13}
% FIRST_POSTED: 2026-10-02
% ENTRY_KIND: statement_only
% UnsolvedMath ID 2302082; source catalogue code AMR-022-2082
% !TeX program = lualatex
% Definition and sources only; this file is not an LLM solution attempt.
\documentclass[11pt,a4paper]{article}
\usepackage[margin=25mm]{geometry}
\usepackage{fontspec}
\setmainfont{lmroman10-regular.otf}[BoldFont=lmroman10-bold.otf,
 ItalicFont=lmroman10-italic.otf,BoldItalicFont=lmroman10-bolditalic.otf]
\usepackage{amsmath,amssymb,mathtools}
\usepackage{unicode-math}
\setmathfont{latinmodern-math.otf}
\AtBeginDocument{\let\setminus\smallsetminus}
\usepackage{xurl,hyperref}
\hypersetup{colorlinks=true,urlcolor=blue}
\setcounter{secnumdepth}{0}
\setlength{\parindent}{0pt}
\setlength{\parskip}{5pt}
\setlength{\emergencystretch}{3em}
\begin{document}
\section*{Research Problems in Function Theory — Problem 2.82}\label{2302082}
{\small UnsolvedMath ID 2302082 · AMR-022-2082 · Analysis · AMR Open Problem Lists}
\subsection{Definitions and mathematical statement}
For an integer $d\ge2$, let $\operatorname{Rat}_d$ be the space of degree-$d$ rational maps of the Riemann sphere. Represent a map by coprime homogeneous degree-$d$ polynomials, modulo common nonzero scaling; the induced coefficient topology is used. A critical point has local mapping degree greater than one. The Julia set $J(g)$ is the complement of the largest open set on which the iterates $\{g^n:n\ge0\}$ form a normal family in the spherical metric.
A point $c$ is preperiodic for $g$ if its forward orbit
$\{c,g(c),g^2(c),\ldots\}$ is finite. It is periodic if $g^m(c)=c$ for some $m\ge1$. Let
\[
\mathcal L_d=\{g\in\operatorname{Rat}_d:
\text{every critical point is preperiodic and none is periodic}\}.
\]
Multiplicity of a critical point does not change the orbit condition; the critical point at infinity is included when present.
Herman's approximation question asks whether
\[
J(f)=\widehat{\mathbb C}\quad\Longrightarrow\quad
f\in\overline{\mathcal L_d}
\qquad(f\in\operatorname{Rat}_d).
\]
Thus, given any such $f$ and any neighbourhood in the fixed-degree coefficient space, one seeks a nearby rational map whose critical orbits are all finite but avoid periodic critical points. The approximants themselves need not have the same critical-point multiplicities as $f$. The closure is taken among degree-$d$ maps, rather than in a compactification allowing common factors and a drop of degree. The assertion is stronger than merely finding maps with many, but not all, preperiodic critical points.
\subsection{Short English statement}
Can every rational map whose Julia set is the whole sphere be approximated by maps with strictly preperiodic critical points?
\subsection{Sources}
\begin{itemize}
\item[S1] ULAM.AI and the UnsolvedMath Contributors, supplied problems.json, record 2302082 (AMR-022-2082), AMR Open Problem Lists. Catalogue identity and source transcription; CC BY 4.0.
\url{https://huggingface.co/datasets/ulamai/UnsolvedMath}.
\item[S2] Hayman - Research Problems in Function Theory (2018), Problem 2.82. Original source indexed by AMR-022-2082.
\url{https://arxiv.org/abs/1809.07200}.
\end{itemize}
\end{document}
